% Paper 1, §2 — Pre-2017 prequel and the 2017 canonical
% Sets up the load-bearing notation, equations, and tensor shapes that
% §3-§12 build on. Math is transcribed verbatim from kb/excerpts/vaswani2017.md;
% pre-2017 lineage is sketched with deepen-cite markers where the KB does not
% yet have primary excerpts (Bahdanau 2014, Gehring 2017, Hochreiter 1997).

\section{Pre-2017 prequel and the 2017 canonical}
\label{sec:prequel-and-canonical}

The modern Transformer is a small set of choices made on top of a
2017 canonical that was itself a tightly-engineered response to the
known pathologies of the recurrent and convolutional sequence-to-sequence
architectures it replaced. \textbf{The thesis of this section: every
axis of variation walked in the rest of this paper (KV-sharing,
attention-implementation, FFN, normalization, positional encoding,
multi-token output, multimodality) is a localized perturbation of the
six-block encoder-decoder of \citet{vaswani2017}.} Establishing what
that 2017 baseline actually computes -- with every variable defined,
every shape labeled, and the four equations on which $\S 3$--$\S 12$
will pivot transcribed verbatim -- is the load-bearing job of this
section. We start with what the Transformer replaced, why it
replaced them, then state the 2017 architecture in the notation we
will carry through the paper.

\subsection{What the Transformer replaced}
\label{sec:prequel}

By the end of 2016 the dominant family for sequence transduction was
the recurrent encoder-decoder with attention. Three pathologies were
known well enough to be the abstract of \citet{vaswani2017} itself
(``dispensing with recurrence and convolutions
entirely'')%
\footnote{KB excerpt: \texttt{kb/excerpts/vaswani2017.md\#abstract}.}
\citep[abstract]{vaswani2017}.

\paragraph{RNN/LSTM seq2seq: the sequential bottleneck.} The
recurrent encoder consumes one token at a time, computing
$\mathbf{h}_t = f(\mathbf{h}_{t-1}, \mathbf{x}_t)$ for $t = 1, \ldots, S$,
which makes the per-step depth equal to the sequence length and
forecloses parallelism along $S$. Long-range gradient flow through
this recurrence vanishes (or, less commonly, explodes) at a rate set
by the operator norm of the Jacobian
$\partial \mathbf{h}_t / \partial \mathbf{h}_{t-1}$; the LSTM cell
\deepencite{hochreiter1997-lstm §3 (gating equations: input, forget,
output gates plus cell-state recurrence)} replaces a saturating
non-linearity with a gated linear cell-state pathway
$\mathbf{c}_t = \mathbf{f}_t \odot \mathbf{c}_{t-1} + \mathbf{i}_t \odot \tilde{\mathbf{c}}_t$
to soften the decay, but does not remove it. Equally damaging at
training time, the $O(S)$ wall-clock step count multiplies with the
depth-$L$ stack so the encoder's wall-clock training scales as
$O(L \cdot S)$ even with full hardware parallelism along the batch
axis. \citet{vaswani2017}~\S 1 names this directly: ``This
inherently sequential nature precludes parallelization within
training examples, which becomes critical at longer sequence lengths
\dots''

\paragraph{Encoder-decoder attention: the right idea, wrong base.}
\citet{vaswani2017}'s ``attention'' did not appear from nothing.
\deepencite{bahdanau2014-attention §A.2 (additive attention scoring
function $a(\mathbf{s}_{i-1}, \mathbf{h}_j) = \mathbf{v}_a^\top
\tanh(W_a \mathbf{s}_{i-1} + U_a \mathbf{h}_j)$ and the soft-alignment
context $\mathbf{c}_i = \sum_j \alpha_{ij} \mathbf{h}_j$)} introduced
the mechanism: at each decoder step $i$, learn weights
$\alpha_{ij} = \softmax_j(a(\mathbf{s}_{i-1}, \mathbf{h}_j))$ over the
encoder's hidden states $\{\mathbf{h}_j\}_{j=1}^{S_{\text{src}}}$, and
read out $\mathbf{c}_i = \sum_j \alpha_{ij} \mathbf{h}_j$ as a
content-addressed summary of the source. The attention idea was sound
and survived intact; what \citet{vaswani2017} would discard is the
\emph{recurrent base} that the 2014--2016 encoder-decoders bolted
attention onto, and the \emph{additive} compatibility function in
favor of a scaled dot-product (\S\ref{sec:canonical-attention} below).

\paragraph{ConvSeq2Seq: parallel but local.} A natural alternative
to recurrence is convolution. \deepencite{gehring2017-convseq2seq §2
(GLU-gated 1D convolutions over source and target with multi-step
attention)} replaces both encoder and decoder recurrences with stacks
of GLU-gated 1D convolutions, gaining within-example parallelism
along $S$ at the price of a $O(\log_k S)$-deep receptive field for
kernel size $k$: long-range dependencies require either deep stacks
or attention bolted on top. This was the strongest direct rival to
the Transformer at the moment of submission and forms the headline
comparison in \citet[Table~2]{vaswani2017}, where the Transformer
matches or exceeds ConvS2S on WMT'14 EN--DE BLEU at a fraction of the
training FLOPs.

\begin{intuition}
Three failure modes, one root cause: each of the pre-2017 designs
forces inter-position information transfer through a structurally
constrained pipe. RNNs route everything through a $D$-dim recurrent
state at $O(S)$ sequential depth; CNNs route through a fixed kernel
of $O(\log_k S)$ depth. Attention discards the pipe. Returning to
math: the canonical attention function (Eq.~\eqref{eq:vaswani-attn}
below) computes an $S \times S$ score matrix in one parallel step,
making every position directly addressable from every other position
at $O(1)$ ``hops'' and $O(S^2)$ FLOPs. The trade is depth-for-quadratic,
and -- as the rest of this paper documents -- the field paid the
quadratic price and built the engineering apparatus to absorb it.
\end{intuition}

\subsection{The 2017 canonical: encoder-decoder structure}
\label{sec:canonical-structure}

The Transformer of \citet{vaswani2017} is a stack of $L$ identical
blocks operating on sequences of $D$-dimensional vectors, with no
recurrence and no convolution. We carry the tensor-shape conventions
of \cref{tab:tensor-shapes} through this section and the rest of the
paper.

\begin{table}[h]
\centering
\renewcommand{\arraystretch}{1.2}
\begin{tabular}{lll}
\toprule
\textbf{Symbol} & \textbf{Meaning} & \textbf{Vaswani-base value} \\
\midrule
$\mathsf{B}$        & batch size                                  & --- (training-dependent) \\
$\mathsf{S}$        & sequence length (tokens)                    & up to 512 (training); -- at inference \\
$\mathsf{D}$        & residual-stream / model dimension           & $512$ \\
$\mathsf{H}$        & number of attention heads                   & $8$ \\
$d_h$               & per-head dimension, $d_h = \mathsf{D}/\mathsf{H}$ & $64$ \\
$d_{\text{ff}}$     & FFN inner dimension                         & $2048$ ($= 4\mathsf{D}$) \\
$\mathsf{V}$        & tokenizer vocabulary size $|V|$             & $\sim$37k (BPE, EN--DE)~\citep[\S 5.1]{vaswani2017} \\
$L$                 & blocks per stack (encoder; decoder)         & $6$ each \\
\bottomrule
\end{tabular}
\caption{Tensor-shape and capacity notation used throughout this
paper, with the Vaswani-base values from \citet[\S 3.1, \S 5.1]{vaswani2017}.
The residual stream carries shape $\BSH$; after the head-split inside
attention it carries $\BHSdh$. Subsequent sections (\S\ref{sec:tokenization-embedding}
onward) hold this notation fixed and label any divergence
(MLA's $d_c$, MoE's $E$ and $k$, GQA's $G$).}
\label{tab:tensor-shapes}
\end{table}

The encoder is a stack of $L = 6$ identical layers; each has two
sublayers, a multi-head self-attention block and a position-wise
feed-forward block, with a residual connection and layer
normalization wrapping each sublayer%
\footnote{KB excerpt: \texttt{kb/excerpts/vaswani2017.md\#sec-3-1}.}.
The decoder also has $L = 6$ layers but inserts a third sublayer --
encoder-decoder cross-attention -- between the self-attention and
the FFN, and masks the self-attention so that position $i$ can attend
only to positions $\leq i$ \citep[\S 3.1]{vaswani2017}. Writing the
sublayer abstractly as $F$ and its input as $\mathbf{x}$, the
canonical placement is what \citet{vaswani2017} state verbatim%
\footnote{KB excerpt: \texttt{kb/excerpts/vaswani2017.md\#sec-3-1}, the same anchor as immediately above.}:
\begin{equation}
  \mathbf{y} \;=\; \mathrm{LayerNorm}\bigl(\mathbf{x} + F(\mathbf{x})\bigr)
  \tag{Vaswani--\S 3.1}
  \label{eq:post-ln-original}
\end{equation}
\citep[\S 3.1]{vaswani2017}. This is the \emph{Post-LN} placement
(norm \emph{after} the residual add), and it is the placement
\citet{radford2019-gpt2} will discard in GPT-2 in favor of Pre-LN
(\S\ref{sec:normalization-axis}); the choice has nontrivial
consequences for training stability at depth \citep{xiong2020}, but
in the 2017 canonical we keep Eq.~\eqref{eq:post-ln-original}
exactly. All sublayers and the embedding layers produce outputs of
the same dimension $\mathsf{D} = 512$, ``to facilitate these residual
connections''~\citep[\S 3.1]{vaswani2017}.

\subsection{Scaled dot-product attention}
\label{sec:canonical-attention}

The attention sublayer is the load-bearing departure from the
recurrent and convolutional baselines. \citet{vaswani2017} define
scaled dot-product attention on per-head matrices
$Q, K \in \R^{\mathsf{S} \times d_k}$, $V \in \R^{\mathsf{S} \times d_v}$
as%
\footnote{KB excerpt: \texttt{kb/excerpts/vaswani2017.md\#sec-3-2-1}.}
\begin{equation}
  \mathrm{Attention}(Q, K, V) \;=\;
  \softmax\!\left(\frac{Q K^\top}{\sqrt{d_k}}\right) V
  \tag{Vaswani--Eq.~1}
  \label{eq:vaswani-attn}
\end{equation}
\citep[\S 3.2.1, Eq.~1]{vaswani2017}. The inner $Q K^\top$ is the
$\mathsf{S} \times \mathsf{S}$ score matrix; row $i$ of the post-softmax
result is a probability distribution over keys, and the multiplication
by $V$ on the right reads out a convex combination of value rows --
\emph{content-based retrieval} from a set of key-value pairs. The
$\sqrt{d_k}$ scaling keeps the pre-softmax logit variance independent
of $d_k$ at initialization: for $q, k$ with i.i.d.\ unit-variance
components, $\mathrm{Var}(q \cdot k) = d_k$~\citep[\S 3.2.1, fn.~4]{vaswani2017}, and
unscaled dot products grow with $d_k$ ``pushing the softmax function
into regions where it has extremely small gradients''~\citep[\S
3.2.1]{vaswani2017}.

\begin{analogy}
Eq.~\eqref{eq:vaswani-attn} is a soft, learnable, parallel
associative-memory lookup. The keys are addresses; the values are
contents; the query is what we are asking. The softmax over
$Q K^\top / \sqrt{d_k}$ is the soft index that, were it a
one-hot vector, would pick a single $(K_j, V_j)$ pair; in the
soft case it is a content-weighted average. Returning to canonical
math: the row $i$ output is
$\sum_j \softmax_j(q_i^\top k_j / \sqrt{d_k}) \, v_j$,
which is precisely a weighted-sum read of $V$ keyed by the
content alignment between $q_i$ and the $k_j$'s.
\end{analogy}

\paragraph{Multi-head.} Rather than running a single attention
function on $D$-dimensional queries, keys, and values,
\citet{vaswani2017} project them $\mathsf{H}$ times to $d_k$, $d_k$,
$d_v$ dimensions and run Eq.~\eqref{eq:vaswani-attn} in
parallel%
\footnote{KB excerpt: \texttt{kb/excerpts/vaswani2017.md\#sec-3-2-2}.}:
\begin{align}
  \mathrm{MultiHead}(Q, K, V)
   &= \mathrm{Concat}(\mathrm{head}_1, \ldots, \mathrm{head}_{\mathsf{H}})\, W^O,
   \notag \\
  \mathrm{head}_i
   &= \mathrm{Attention}(Q W_i^Q, K W_i^K, V W_i^V),
   \tag{Vaswani--Eq.~2--3}
  \label{eq:vaswani-multihead}
\end{align}
\citep[\S 3.2.2, Eq.~2--3]{vaswani2017}, with projection matrices
$W_i^Q, W_i^K \in \R^{\mathsf{D} \times d_k}$,
$W_i^V \in \R^{\mathsf{D} \times d_v}$, and
$W^O \in \R^{\mathsf{H} d_v \times \mathsf{D}}$. The 2017 base
chooses $\mathsf{H} = 8$ heads and
$d_k = d_v = d_h = \mathsf{D}/\mathsf{H} = 64$, so the per-head
projections \emph{partition} the model dimension across heads and
the total compute matches a single-head full-width attention
\citep[\S 3.2.2]{vaswani2017}. ``Multi-head attention allows the
model to jointly attend to information from different
representation subspaces at different positions; with a single
attention head, averaging inhibits this''~\citep[\S
3.2.2]{vaswani2017}\footnote{KB excerpt: \texttt{kb/excerpts/vaswani2017.md\#sec-3-2-2-motivation}.}.
After projection the attention tensors carry shape $\BHSdh$; the
score matrix has shape
$\mathsf{B}\!\times\!\mathsf{H}\!\times\!\mathsf{S}\!\times\!\mathsf{S}$;
the post-softmax-times-$V$ output returns to $\BHSdh$ and the
$W^O$ projection collapses heads back to $\BSH$. The shape
$\BHSdh$ is the lens the rest of this paper is read through.

\paragraph{Three uses of attention.} The 2017 model uses
multi-head attention in three places~\citep[\S
3.2.3]{vaswani2017}\footnote{KB excerpt: \texttt{kb/excerpts/vaswani2017.md\#sec-3-2-3}.}:
encoder self-attention (queries, keys, and values all from the
previous encoder layer), decoder \emph{masked} self-attention (same
sources, but with positions $> i$ masked to $-\infty$ inside the
softmax to preserve autoregression), and encoder-decoder
cross-attention (queries from the previous decoder layer; keys and
values from the encoder output). The decoder-only LLMs of $\S 3$
onward keep only the second of these three -- masked self-attention
-- and the resulting single causal stack is what survives into
2026.

\subsection{Position-wise feed-forward network}
\label{sec:canonical-ffn}

Each block contains a position-wise feed-forward sublayer applied
identically and independently to each position%
\footnote{KB excerpt: \texttt{kb/excerpts/vaswani2017.md\#sec-3-3}.}:
\begin{equation}
  \mathrm{FFN}(\mathbf{x}) \;=\; \max(0, \mathbf{x} W_1 + b_1) \, W_2 + b_2
  \tag{Vaswani--Eq.~2 (FFN)}
  \label{eq:vaswani-ffn}
\end{equation}
\citep[\S 3.3]{vaswani2017}, with
$W_1 \in \R^{\mathsf{D} \times d_{\text{ff}}}$,
$W_2 \in \R^{d_{\text{ff}} \times \mathsf{D}}$, and the canonical
choice $d_{\text{ff}} = 4 \mathsf{D}$ (so $d_{\text{ff}} = 2048$ in
the base configuration). \citet{vaswani2017} note this is also
``two convolutions with kernel size 1''~\citep[\S 3.3]{vaswani2017}
-- the FFN does not mix across positions, only across channels at
each position. The same $(W_1, b_1, W_2, b_2)$ is shared across all
$\mathsf{S}$ positions but \emph{not} across the $L$ layers; each
layer learns its own pair. The activation is ReLU
($\max(0, \cdot)$). The factorization of every block into ``mix
across positions, then mix across channels'' is the structural
property the rest of the paper preserves: $\S 5$ replaces the
attention sublayer's KV-sharing pattern, $\S 7$ replaces the FFN's
ReLU-2-matrix form (with SwiGLU and then MoE), and the strict
position/channel split survives every variant.

\subsection{Sinusoidal positional encoding}
\label{sec:canonical-position}

Because the model contains no recurrence or convolution, position
information must be injected explicitly%
\footnote{KB excerpt: \texttt{kb/excerpts/vaswani2017.md\#sec-3-5}.}.
\citet{vaswani2017} add a fixed sinusoidal encoding to the input
embeddings at the bottom of each stack:
\begin{equation}
  PE_{(pos, 2i)}   = \sin(pos / 10000^{2i/\mathsf{D}}),
  \qquad
  PE_{(pos, 2i+1)} = \cos(pos / 10000^{2i/\mathsf{D}}),
  \tag{Vaswani--PE}
  \label{eq:vaswani-pe}
\end{equation}
\citep[\S 3.5]{vaswani2017}, where $pos$ is the integer position
and $i$ indexes the embedding dimension. The two halves of the
embedding are filled with a sinusoid of geometrically progressing
wavelengths from $2\pi$ to $10000 \cdot 2\pi$, and
\citet{vaswani2017} motivate the choice as making relative position
linearly recoverable: ``for any fixed offset $k$, $PE_{pos + k}$
can be represented as a linear function of $PE_{pos}$''~\citep[\S
3.5]{vaswani2017}. They also report that learned absolute embeddings
match in quality, and pick the sinusoidal version for the
extrapolation-to-longer-sequences argument it appears (in 2017) to
license. The lineage that follows -- learned absolute, T5 relative
bias, ALiBi, RoPE -- is the subject of \S\ref{sec:positional-encoding};
RoPE in particular makes a different bet on \emph{multiplicative}
position injection inside the attention dot product, and is the
2026 default.

\subsection{Embeddings, weight tying, and softmax output}
\label{sec:canonical-embed}

The token-ID-to-vector and vector-to-logit interfaces are also
specified in the paper%
\footnote{KB excerpt: \texttt{kb/excerpts/vaswani2017.md\#sec-3-4}.}.
The input embedding $E_{\text{in}} \in \R^{\mathsf{V} \times \mathsf{D}}$
maps token IDs to $\mathsf{D}$-vectors; the LM head
$E_{\text{out}} \in \R^{\mathsf{D} \times \mathsf{V}}$ projects the
final residual stream to vocabulary logits; a softmax produces
next-token probabilities. \citet{vaswani2017} adopt the
\emph{weight-tying} scheme of \citet{press2017-tied-embed}: the
matrix used as $E_{\text{in}}$ is also used (transposed) as
$E_{\text{out}}$, so a single $\mathsf{V} \times \mathsf{D}$ table
serves both interfaces, and ``in the embedding layers, we multiply
those weights by $\sqrt{\mathsf{D}}$''~\citep[\S 3.4]{vaswani2017}. In
the notation of $\S 3$ this is
\begin{equation}
  E_{\text{out}} \;=\; E_{\text{in}}^{\!\top},
  \tag{Press--Vaswani tying}
  \label{eq:weight-tying}
\end{equation}
\citep{press2017-tied-embed}. Tying saves $\mathsf{V} \times \mathsf{D}$
parameters -- a substantial fraction of the 2017 base model's $\sim 65$M
parameters when $\mathsf{V} \approx 37\text{k}$ and $\mathsf{D} = 512$.
\S\ref{sec:tokenization-embedding} discusses why frontier-2026
\emph{untied} the I/O embeddings as $\mathsf{V}$ and $\mathsf{D}$ both
grew (the relative parameter share of the embedding table fell, and
the inductive bias that input and output token spaces should share
geometry weakened); GPT-2 \citep[\S 2.3]{radford2019-gpt2} kept
tying, but LLaMA-3 and DeepSeek-V3 do not.

\subsection{The block in one picture, and what comes next}
\label{sec:canonical-block-picture}

Combining Eqs.~\eqref{eq:post-ln-original}, \eqref{eq:vaswani-multihead},
and \eqref{eq:vaswani-ffn}, one decoder block in the 2017 canonical
takes residual stream $X^\ell \in \R^{\mathsf{B} \times \mathsf{S} \times \mathsf{D}}$
to $X^{\ell+1}$ via
\begin{align}
  X^{\ell + 1/2} &= \mathrm{LayerNorm}\!\bigl(X^\ell + \mathrm{MaskedMHA}(X^\ell)\bigr),
  \\
  X^{\ell + 1}   &= \mathrm{LayerNorm}\!\bigl(X^{\ell + 1/2} + \mathrm{FFN}(X^{\ell + 1/2})\bigr),
\end{align}
with positional information added once at the bottom via
Eq.~\eqref{eq:vaswani-pe}, the embedding via
Eq.~\eqref{eq:weight-tying}, and the LM head re-using
$E_{\text{in}}^\top$. The full hyperparameter row of the base
model is $L = 6$, $\mathsf{D} = 512$, $\mathsf{H} = 8$,
$d_h = 64$, $d_{\text{ff}} = 2048$, $\mathsf{V} \approx 37\text{k}$
\citep[\S 3.1, \S 5.1]{vaswani2017}.

\paragraph{What changes after 2017, by axis.} Every subsequent
section of this paper changes \emph{exactly one} of the components
just defined, holding the others fixed. \S\ref{sec:tokenization-embedding}
takes the tokenization, the input embedding, and the unembedding,
and walks subword $\to$ byte-level $\to$ tokenizer-free, plus the
tied-vs-untied decision of Eq.~\eqref{eq:weight-tying}.
\S\ref{sec:positional-encoding} takes the sinusoidal
Eq.~\eqref{eq:vaswani-pe} and walks learned $\to$ T5 relative bias
$\to$ ALiBi $\to$ RoPE $\to$ NoPE.
\S\ref{sec:attention-kv-sharing} keeps Eq.~\eqref{eq:vaswani-attn}
unchanged but varies how the $\mathsf{H}$ heads in
Eq.~\eqref{eq:vaswani-multihead} share their $K$ and $V$
projections (MHA $\to$ MQA $\to$ GQA $\to$ MLA).
\S\ref{sec:attention-mechanism-axis} keeps both Eqs.~\eqref{eq:vaswani-attn}
and \eqref{eq:vaswani-multihead} mathematically unchanged but
schedules the same computation differently onto GPU memory hierarchy
(FA1 $\to$ FA2 $\to$ FA3 $\to$ NSA).
\S\ref{sec:feedforward-axis} replaces the ReLU-2-matrix
Eq.~\eqref{eq:vaswani-ffn} with SwiGLU and then with MoE.
\S\ref{sec:normalization-axis} replaces LayerNorm in
Eq.~\eqref{eq:post-ln-original} with RMSNorm and moves it to
Pre-LN.
\S\ref{sec:multi-token-output} adds heads to the LM head of
\S\ref{sec:canonical-embed}.
\S\ref{sec:multimodal-extensions} extends the embedding interface
of \S\ref{sec:canonical-embed} to non-text inputs.

That \emph{everything else} -- the residual structure, the
position/channel factorization, the softmax-over-keys, the head
partition, the $\mathsf{B}\!\times\!\mathsf{S}\!\times\!\mathsf{D}$
residual stream as the workspace -- survives across all eight axes
is the empirical fact this paper is built around. We turn first to
the discrete-token interface of \S\ref{sec:tokenization-embedding}.
